% SUBSECTION 2.4 "Methodology and Empirical Design" -- DRAFT for advisor, then user (PDF). DASH-FREE (no '---'/'--').
% Markers %%Pn%%. Equations MA1/MA2/MA3 VERIFIED against code this session:
%   MA1 gen_empire_did_table.py:141 (CTRL:47, PreAnnounceQtr=1[e=-1]:127, post dropped:126, firm-cluster:143)
%   MA2 empire_drop_test.py:149 (BINS PRE2/PRE1/GAP/POST :137-140, GAP=announced&~closed, cap+4:133, Wald PRE1-GAP:172)
%   MA3 empire_cashspec_interaction.py:90 (PreAnn_cash/stock:71-72, sample cq<dq_first:69, Wald cash-stock:100-106, cash lag:84)
% SEAM-E bridge: beta=>theta/kappa (2.2). NO NUMBERS (design only; C1/C2/C6 framing). NEW cites opler/bates/petersen/cameron
%   DROPPED (unverified, not in bib) -> generic 'standard cash-determinants regression' + harford callback (RULE-COH-3); flag to NLM-add later.
% Phase C self-scan: no '---'/'--'; RULE-COH-2 (descriptive, no parallel-trends/causal); design->hypothesis 1:1; two control sets distinct;
%   GAP discriminator = uncertainty drop (cash trace is the other clock); selection + generated-regressand disclosed; no overclaim.

%%P1%%
Each hypothesis of Section~2.2 maps to an estimating equation, and all three share a common panel frame: a within-firm, two-way fixed-effects regression of an outcome on event-time treatment indicators, estimated on quarterly firm observations with standard errors clustered by firm. The main design tests H1. Writing $Y_{it}$ for the outcome of firm $i$ in quarter $t$,
\begin{equation}
Y_{it} = \beta\,\mathrm{PreAnnounceQtr}_{it} + \gamma' X_{it} + \phi_i + \tau_t + \varepsilon_{it},
\end{equation}
where $X_{it}$ is a vector of firm controls (leverage, log assets, Tobin's $Q$, return on assets, capital expenditure, a dividend indicator, and scaled cash flow from operations), $\phi_i$ is a firm fixed effect, and $\tau_t$ is a calendar year-quarter fixed effect. The treatment $\mathrm{PreAnnounceQtr}_{it}=\mathbf{1}[e=-1]$ marks the call in the quarter before the firm's first cash acquisition; post-announcement quarters ($e\ge 0$) are dropped, so the design reads the run-up and not the aftermath, while firms with no cash acquisition, for which $e$ is undefined, anchor the calendar-time effects. The coefficient $\beta$ is the estimand: it estimates the pre-announcement shift $\theta_{-1}$ of Section~2.2, a within-firm contrast between a treated firm's pre-announcement quarter and its own ordinary quarters. We read it descriptively, as an anticipation contrast, and make no causal or parallel-trends claim.

%%P2%%
H1b requires tracing the outcome across event time rather than at a single point, so the main design generalizes to an event study that replaces the single indicator with a set of event-time bins:
\begin{equation}
Y_{it} = \sum_{b}\beta_b\,\mathrm{Bin}_{b,it} + \gamma' X_{it} + \phi_i + \tau_t + \varepsilon_{it},
\end{equation}
with four bins: $\mathrm{PRE2}$ ($e=-2$), $\mathrm{PRE1}$ ($e=-1$), $\mathrm{GAP}$ (announced but not yet completed), and $\mathrm{POST}$ (completed). The omitted baseline is the set of quarters at $e\le -3$ together with never-acquirers; the post-announcement window is capped at four quarters and truncated at a firm's next announcement. The discriminating window is the gap. If residual uncertainty has resolved by announcement, its coefficient should fall to zero at $\mathrm{GAP}$, when the deal is public but the cash has not yet left, so a Wald test of $\beta_{\mathrm{PRE1}}-\beta_{\mathrm{GAP}}>0$ carries the empirical content of H1b; the same specification, estimated for the cash ratio in place of uncertainty, traces the persistence of the cash on the other clock, a persistence that is largely mechanical, since the cash is paid only at completion and so sits on the books through the gap. $\mathrm{PRE2}$ is a pre-trend check, expected to be near zero. The bins are estimated for cash acquirers and, as a placebo, for stock acquirers, and every coefficient and Wald contrast is reported two-tailed.

%%P3%%
H1a is a claim about the difference between cash and stock acquirers, so it is tested in a single pooled specification carrying both treatments at once:
\begin{equation}
Y_{it} = \beta_c\,\mathrm{PreAnnCash}_{it} + \beta_s\,\mathrm{PreAnnStock}_{it} + \gamma' X_{it} + \phi_i + \tau_t + \varepsilon_{it},
\end{equation}
where $\mathrm{PreAnnCash}_{it}$ and $\mathrm{PreAnnStock}_{it}$ mark the quarter before a firm's first cash and first stock acquisition. The sample keeps only the quarters before a firm's first deal of either type, plus firms that make neither, which form the fixed-effects baseline. The formal cash-specificity test is the single linear restriction $\beta_c-\beta_s>0$, evaluated by a Wald test; this difference, and not a comparison of one significant coefficient against one insignificant one, is what estimates $\theta_{-1}^{\,\mathrm{cash}}-\theta_{-1}^{\,\mathrm{stock}}$ in H1a. When the outcome is the cash ratio rather than residual uncertainty, the specification adds the cash ratio's own one-quarter lag, so the pre-announcement coefficient tracks the change in cash rather than its level; the residual outcome takes no such lag.

%%P4%%
The identifying assumptions are descriptive and common to the three designs. Identification is within firm, through the firm fixed effects, and within calendar time, through the year-quarter fixed effects; the estimand in each case is a within-firm pre-announcement mean shift, measured against a treated firm's own ordinary quarters; firms that never acquire anchor the calendar-time effects rather than the within-firm contrast itself. The $\mathrm{PRE2}$ bin is a validity device rather than an identification proof: a pre-trend near zero is consistent with the reading, but we claim no parallel-trends result and no causal effect. The mapping from design to hypothesis is one to one: the single-indicator regression tests H1, the event study tests H1b, and the pooled interacted regression tests H1a. The control pool itself differs across designs: in the single-indicator and event-study regressions a firm remains in the baseline until its first cash acquisition, so stock-only acquirers sit among the controls, whereas the pooled regression restricts the baseline to firms that make neither a cash nor a stock acquisition. Two distinct control layers operate and should not be conflated. The residual outcome is already net of a first set of controls, the linguistic and performance variables of the decomposition in Section~2.3, estimated with CEO and year fixed effects; the designs here add a second set, the firm-financial controls $X_{it}$, with firm and finer year-quarter fixed effects. The two layers serve different purposes and are estimated at different stages.

%%P5%%
Several choices that do not belong to a single equation are recorded here. On inference, standard errors are clustered by firm throughout, while some validity tables cluster two ways, by firm and by calendar quarter; focal treatment tests are reported two-tailed. On contamination, deal timing is endogenous, and because only a firm's first deal sets the event clock, later deals can pollute the baseline; we mitigate this by truncating each firm's post-window at its second announcement, but the never-acquirer baseline and the first-deal focus remain assumptions, not identification. On functional form, the outcomes are levels or word shares rather than counts, the cash-ratio regressions carry a partial-adjustment lag because cash is sticky whereas the residual regressions do not, and, because the residual is a generated regressand from Section~2.3, the two-step standard-error caveat carries over to every design here. On selection, the residual outcome exists only for CEOs with at least five calls, in the non-financial, non-utility sample where the decomposition is run, so the uncertainty regressions describe a larger, more-covered set of firms than the cash-ratio regressions, a limit on external validity we disclose. The firm-financial controls follow the standard cash-determinants regression, read mechanically, with cash treated as a by-product of retained free cash flow rather than a war chest built for a deal.
%%END%%
